% Section VII (plan v2): certificates that bound events. Drafted 2026-09-06; subsections A-C are added as their results arrive (feat-019/020/021); D uses results/warped_anchor.csv (feat-022).
\section{Certificates That Bound Events}\label{sec:certificates}

The gaps found so far are properties of the guarantee, not of the implementation, and each admits a repair inside the decoder. This section implements three and prices them: a pathwise budget on the realised log-ratio, which certifies every event (Section~\ref{sec:pathwise}); a concentration certificate for the unchanged KL decoder, computed from quantities the solver already evaluates (Section~\ref{sec:concentration}); and a per-user budget with a cap on the bank (Section~\ref{sec:odometer}). Section~\ref{sec:warped} then shows that all of them, and Theorem~3.1 itself, are stated relative to an anchor that the deployer's decoding settings modify.

\subsection{A pathwise budget certifies events}\label{sec:pathwise}
The KL certificate bounds an expectation because the decoder charges one: $a_t = D_{\mathrm{KL}}(p_\theta \,\|\, p_s)$ averages the log-ratio over tokens that were not sampled. Charging instead the log-ratio of the sampled token, $r_t = \log p_\theta(y_t) - \log p_s(y_t)$, gives a decoder whose invariant is the realised ratio itself. At every step it takes the largest $\theta_t \in [0, 1]$ whose maximal log-ratio $m_t(\theta) = \max_v \log p_\theta(v)/p_s(v)$ fits the allowance $k^{\mathrm{pw}}_t = \max\{0, (t+1)k - \delta_{\mathrm{init}}(x) - R_{t-1}\}$, with $R_{t-1} = \sum_{i < t} r_i$; $m_t$ increases in $\theta$ from $m_t(0) = 0$, so a bisection finds $\theta_t$, and $\theta_t = 1$ whenever the risky model fits. Proposition~\ref{prop:pathwise} shows that every output then satisfies $L(y) \le K - \delta_{\mathrm{init}}(x)$, the $\Delta_{\max}$ form of near-access-freeness, so $p_\theta(E \mid x) \le e^{K} p_s(E \mid x)$ for every event $E$ by Lemma~2.2 of Vyas et al., with the KL form as a corollary. The banking rule is unchanged, a token the anchor preferred refunds allowance, and the per-path sum of KL spends is no longer bounded and need not be (it reaches $1.12K$ at $k=1$ while its expectation stays below $K$).

We ran the composition attack against the pathwise decoder on the same passages, seeds, and budgets. The invariant held on all $23{,}844$ budgeted queries against the memorising 8B model and on the $1{,}250$ queries of the 70B audit at the authors' settings, with no overshoot. The event bound costs the adversary a little at every budget: with the memorising model, oracle windows recover $0.06$, $0.19$, $0.39$, and $0.73$ of a passage at $k = 3$, $5$, $10$, and $20$ against $0.10$, $0.23$, $0.41$, and $0.79$ under the KL decoder, chained windows $0.00$, $0.02$, $0.09$, and $0.44$ against $0.01$, $0.03$, $0.11$, and $0.50$, and a single query $0.45$ against $0.48$ at $k=20$; at $k=50$ both return the unconstrained model. The difference is small because the decoders differ only on the few tokens whose maximal log-ratio exceeds the average by more than the allowance absorbs. What changes is what can be said: the same numbers are now guaranteed for every output, and the reconstruction rate of any window is bounded by $\exp(K_i - S_i)$ (Proposition~\ref{prop:cost}). With the 70B model at the authors' settings the pathwise decoder holds a single query to $0.02$, $0.05$, $0.09$, and $0.31$ and oracle windows to $0.02$, $0.07$, $0.22$, and $0.55$ at $k = 3$, $5$, $10$, and $20$, against $0.31$ and $0.61$ for the sampled model alone.

The price is paid on ordinary prompts. On the sweep of Section~\ref{sec:sweep}, with both decoders on the same $900$ prompts, seeds, and budgets ($2{,}700$ trajectories per budget), the pathwise decoder intervenes more often at every budget: the constraint is active on $56\%$ of decode steps against $43\%$ at $k=0.5$, $14.6\%$ against $5.0\%$ at $k=1$, $1.9\%$ against $0.3\%$ at $k=3$, $1.1\%$ against $0.2\%$ at $k=5$, and $0.5\%$ against $0.01\%$ at $k=10$, and it forces the anchor on $6.1\%$ of steps against $2.2\%$ at $k=3$ (at $k \le 1$ the forced fraction is the prefix debt's and the same for both). The risky model is served unchanged on $92.0\%$ of steps against $97.6\%$ at $k=3$ and on $99.5\%$ against $99.9\%$ at $k=10$; at $k=20$ the two decoders coincide. The risky model's own log-loss on the generated text is within $0.15$ nats per token of the KL decoder's throughout, and all $16{,}200$ pathwise trajectories satisfied $L(y) \le \max\{0, B\}$.

Proposition~\ref{prop:cost} bounds what re-sampling buys the adversary: a window with anchor surprisal $S_i$ and budget $K_i = kL$ is reproduced exactly with probability at most $e^{K_i - S_i}$ per attempt. We measured it with an adversary who re-samples each $50$-token window up to eight times until it matches (50 passages, 297 windows, oracle prefixes). Where the bound is informative it held: at $k=2$ it is below one for $253$ windows and sums to $0.003$ expected reproductions per pass, at $k=3$ for $161$ windows summing to $3.0$, and no window was reproduced exactly at either budget in $2{,}376$ attempts. It is vacuous at the budgets where the attack succeeds, since a window costs the anchor about $150$ nats and $K_i$ is $250$ nats already at $k=5$; there the adversary's cost is set by the budget binding rather than by the certificate, at $3.8$, $38$, and $153$ exact reproductions per thousand attempts at $k=5$, $10$, and $20$ under the pathwise decoder ($9.2$, $44$, and $178$ under the KL decoder, $180$ unconstrained). The bound binds only below about three nats per token for these passages, the regime in which Section~\ref{sec:certificate} found the certificate informative.

\subsection{What the KL certificate concentrates to}\label{sec:concentration}
A deployer who keeps the KL decoder can still say something about single outputs, because the gap between the realised ratio and the charged spend is a martingale whose increments the decoder already computes. Write $M_t = \sum_{i \le t}(r_i - a_i)$; then $L(y) = Z_T + M_T$ with $Z_T \le K$ on every path, so an output exceeds the certificate by $\lambda$ only if the martingale does. Each increment is bounded above by $b_i = m_i - a_i$, the largest log-ratio the served distribution assigns to any token minus the step's spend, and has conditional variance $\sigma_i^2 = \mathrm{Var}_{p_{\theta_i}}[r_i]$, and both are functions of the prefix alone, evaluated before the token is drawn. Proposition~\ref{prop:conc} applies Freedman's inequality~\cite{freedman1975tail,tropp2011freedman}: if $b$ bounds every $b_i$ and $v$ bounds $\sum_i \sigma_i^2$ on every path, then $\Pr[L(y) \ge K + \lambda] \le \delta(\lambda) = \exp(-\lambda^2 / (2(v + b\lambda/3)))$. This is the concentration assumption that Lemma~2.3 of Vyas et al.\ requires for the KL form to bound events, stated in terms the decoder can enforce, by clipping $b_i$ and $\sum_i \sigma_i^2$ as it clips $a_i$, or verify after the fact from the logs. The certificate it yields is weaker than the pathwise one, since it is a bound on the probability of exceeding $K + \lambda$ rather than a bound at $K$, and its strength depends on how large $v$ and $b$ turn out to be at the deployed budget.

We evaluated the certificate on the small-budget sweep of Section~\ref{sec:sweep} at $k=0.5$ and $k=1$ ($2{,}700$ trajectories of $200$ tokens per budget, six prompt classes), taking $v$ and $b$ as the $90$th percentiles of the logged $\sum_i \sigma_i^2$ and $\max_i b_i$ across trajectories, which a deployer could enforce by clipping. At $k=0.5$ ($K=100$) the summed variance has median $217$ (90th percentile $236$) and the largest per-step excess median $14.6$ (90th percentile $19.0$, maximum $32.5$), so the bound gives $\Pr[L(y) \ge K+50] \le 0.11$ and $\Pr[L(y) \ge 2K] \le 0.004$ against empirical frequencies of $0.001$ and $0$; at $k=1$ ($K=200$) it gives $\Pr[L(y) \ge 1.5K] \le 0.009$ against $0$. From $k=3$ to $k=20$ the logged variance saturates at a median of $350$ (90th percentile $450$) and $b$ at $15$ (maximum $29$), since the decoder is nearly the identity and the increments are the risky model's own, so the bound gives $\Pr[L(y) \ge K+100] \le 0.011$ at every one of these budgets, $1.17K$ at $k=3$ and $1.03K$ at $k=20$, against an empirical frequency of $0$ and no output above $K$. The certificate is therefore informative, and loose by two orders of magnitude, because Freedman's inequality pays on every path for a worst-case increment of $15$ to $30$ nats that few tokens produce; it says nothing about the $34\%$ ($k=0.5$) and $8\%$ ($k=1$) of outputs that exceed $K$ itself, since $\delta(0)=1$. Tightening it means bounding $b_i$ per step, which the pathwise decoder does in full.

\subsection{The certificate depends on the decoding settings}\label{sec:warped}
He et al.\ apply the temperature and the repetition penalty to both models' logits before the KL solve (their Appendix~B) and decode books at temperature $0.7$ with penalty $1.1$ (their Appendix~D.1), so Theorem~3.1 certifies proximity to the \emph{warped} anchor $p_s^{(\tau,\rho)}$ and the surprisal in Eq.~\eqref{eq:cap} must be computed under it (Remark~\ref{rem:warp}). For the 758 CopyBench passages, tempering sharpens the anchor and lowers the probability it gives to text that is not its own most likely continuation, so the median surprisal rises from $3.20$ nats per token at temperature $1$ to $3.55$ at the authors' setting, $4.30$ at temperature $0.5$, and $6.53$ at $0.3$ (penalty $1.1$; the penalty alone moves it by at most $6\%$), and the certificate is less often vacuous at small budgets: at $k=1$ the vacuous fraction falls from $44\%$ to $23\%$, $6\%$, and $0\%$, and at $k=1.5$ from $99\%$ to $94\%$, $69\%$, and $7\%$. At $k \ge 3$ it is vacuous for every passage under every setting except temperature $0.3$ ($97\%$). A published $K$ therefore means a different bound under different decoding settings, which the deployer chooses after the theorem is proved, so the settings belong on the certificate; and since a lower temperature also sharpens the served distribution, the attacks at the authors' setting (Section~\ref{sec:natural}) show which effect prevails.



\subsection{What the anchor already knows}\label{sec:leakage}
Near-access-freeness is relative to the anchor: text the anchor assigns high probability to is protected by no $K$, and He et al.\ note that famous fragments of protected works reach openly licensed corpora through quotation (their Appendix~A.3). We measured the leak. For 50 public-domain novels in Project Gutenberg, which the Common Pile includes, the anchor's per-token surprisal of the opening 100 tokens has median $3.11$ nats against $3.33$ for a passage $20{,}000$ characters into the same book (lower for 34 of 50), and the two most quoted openings in English, of \emph{A Tale of Two Cities} and \emph{The Adventures of Sherlock Holmes}, cost $1.2$ nats per token. The leak stops short of reproduction: the anchor alone, prompted with the first sentence of each opening, recovers none of the 50 continuations. For the sixteen CopyBench novels the widely quoted first sentence costs $4.35$ nats per token against $4.73$ for a passage of the same length from the same book (lower for 13 of 16; $2.8$ for \emph{Harry Potter}). This is counterfactual memorisation~\cite{zhang2023counterfactual} of famous lines, which Eq.~\eqref{eq:cap} charges correctly by lowering $S(x)$ only where the text is famous; a certificate card should report the anchor's surprisal on the reference set, since a low value on a work is a fact about the anchor's corpus, not about the mechanism.

\subsection{Accounting across queries}\label{sec:odometer}
Near-access-freeness does not compose~\cite{cohen2025blameless}, and Section~\ref{sec:attack-results} shows the consequence. Differential privacy meets the same problem with a per-user budget tracked by an odometer~\cite{rogers2016odometers,whitehouse2023fully}, and the device transfers because the spend of every query is a number the decoder already computes. We replayed the logs of the composition attack (memorising model, 100 passages) against a per-user budget $B_{\mathrm{user}}$: queries are served in the attack's order until the cumulative spend would exceed the budget, the user is then refused, and reconstruction is scored on the outputs served (Table~\ref{tab:odometer}).

\begin{table}[t]
\centering
\caption{Per-user budget replayed on the composition attack (memorising 8B model, 100 passages, 50-token windows): mean total spend of the uncut attack in nats (median anchor surprisal of a passage: 849), and mean recall of the outputs served before the budget is exhausted, with the fraction of users cut off in parentheses.}
\label{tab:odometer}
\small
\setlength{\tabcolsep}{4pt}
\begin{tabular}{@{}lrrrrr@{}}
\toprule
& & \multicolumn{4}{c}{recall at $B_{\mathrm{user}}$ (users cut)} \\
\cmidrule(l){3-6}
$k$ & spend & 200 & 400 & 800 & $\infty$ \\
\midrule
\multicolumn{6}{@{}l}{\emph{oracle-prefix windows}} \\
3 & 586 & 0.01 (100\%) & 0.06 (98\%) & 0.10 (0\%) & 0.10 \\
5 & 670 & 0.02 (100\%) & 0.11 (100\%) & 0.23 (5\%) & 0.23 \\
10 & 739 & 0.06 (100\%) & 0.18 (100\%) & 0.40 (29\%) & 0.41 \\
20 & 860 & 0.12 (100\%) & 0.30 (100\%) & 0.71 (77\%) & 0.79 \\
\midrule
\multicolumn{6}{@{}l}{\emph{chained windows}} \\
3 & 483 & 0.01 (94\%) & 0.01 (82\%) & 0.01 (0\%) & 0.01 \\
5 & 558 & 0.02 (99\%) & 0.03 (93\%) & 0.03 (0\%) & 0.03 \\
10 & 609 & 0.06 (98\%) & 0.09 (96\%) & 0.11 (3\%) & 0.11 \\
20 & 762 & 0.12 (100\%) & 0.26 (100\%) & 0.47 (41\%) & 0.50 \\
\bottomrule
\end{tabular}
\end{table}

The number that sets the budget is the anchor's surprisal of the work. Reconstructing a passage costs the attack $586$ nats at $k=3$ and $860$ at $k=20$, against a median surprisal of $849$; at $k=20$ the decoder is nearly the identity and every token of the passage is charged in full. A budget equal to the surprisal therefore changes little: at $800$ nats $77\%$ of oracle users are cut off, but the outputs served before the cut already reproduce $71\%$ of the passage. A budget of $400$ nats, half the surprisal, cuts $82$ to $100\%$ of users and bounds reconstruction to $30\%$ (oracle) and $26\%$ (chained); $200$ nats bounds it to $12\%$. The relation is close to proportional: a budget of $B_{\mathrm{user}}$ nats buys about $B_{\mathrm{user}}/3.3$ tokens of a passage at any $k$, and Proposition~\ref{prop:cost} gives the same relation as a lower bound on the adversary's samples once the budget is below the surprisal. The price is that legitimate use is charged at the same rate: a single $260$-token completion spends $576$ to $761$ nats, so $400$ nats refuses every such query and $800$ admits about one per user.

\paragraph{Capping the bank}
The banking rule of Eq.~\eqref{eq:banking} is a token bucket of unbounded depth, and Section~\ref{sec:attack-results} showed that at $k=10$ the bank, rather than the rate, pays for most of a memorised passage. A cap $D$ on the bank discards allowance above $D$, so no $W$ consecutive steps can spend more than $Wk + D$, and $D = k$ is a plain rate limit. We decoded the attack again with $D = k$ and $D = 5k$ on the same seeds ($3{,}804$ queries, all within budget). At $k=10$ the rate limit cuts single-query reconstruction from $0.20$ to $0.12$ and chained windows from $0.11$ to $0.07$, $D = 5k$ gives $0.17$ and $0.09$, and oracle windows stay at $0.39$ to $0.41$ at every depth, because a $50$-token window at $10$ nats per step pays for its tokens without saving. At $k=20$ the cap changes nothing ($0.46$, $0.80$, and $0.47$ against $0.48$, $0.79$, and $0.50$), as the budget-path analysis predicts, since the rate alone exceeds the anchor's surprisal on $92\%$ of passages. The cap bounds bursts, which a deployer may want for its own sake, but not reconstruction.


